% META: {"id":"1NSI-TC-CO-021","chapitre":"1NSI-TYPES-CONSTRUITS","type_objet":"corrige","capacites":["1NSI-TYPES-CONSTRUITS-C2"],"mode_creation":"ex_nihilo","sources_inspiration":[],"status":"approved","exercice_ref":"1NSI-TC-EX-021","origin":"needs_review","status_history":["needs_review","audited","accepted","approved"],"release_acceptance":"RELEASE_OWNER_BATCH_ACCEPTANCE","acceptance_closure_digest":"sha256:9b3ccf9a81c5520fbb7e03b7b2d3e7bf2057a02834b6d908bf3be5f49f3b553f","acceptance_receipt_digest":"sha256:4fad86f0282e0fa4e0419cca1ad6379ae7af5a280c04a4a90880f90a1c044242"}
% BEGIN-VERIFY
% def decaler(tab):
%     dernier = tab[-1]
%     for i in range(len(tab) - 1, 0, -1):
%         tab[i] = tab[i - 1]
%     tab[0] = dernier
% scores = [10, 20, 30, 40]
% decaler(scores)
% assert scores == [40, 10, 20, 30]
% END-VERIFY
\begin{corrige}{1NSI-TC-EX-021}
\begin{enumerate}
  \item Le programme affiche \lstinline{[40, 10, 20, 30]}.

  Détail des itérations (on sauvegarde d'abord \lstinline{dernier = 40}) :
  \begin{itemize}
    \item $i=3$ : \lstinline{tab[3] = tab[2]} $\to$ \lstinline{[10, 20, 30, 30]}
    \item $i=2$ : \lstinline{tab[2] = tab[1]} $\to$ \lstinline{[10, 20, 20, 30]}
    \item $i=1$ : \lstinline{tab[1] = tab[0]} $\to$ \lstinline{[10, 10, 20, 30]}
  \end{itemize}
  Puis \lstinline{tab[0] = dernier} $\to$ \lstinline{[40, 10, 20, 30]}.

  \item La fonction effectue une \textbf{rotation droite} du tableau : le dernier élément passe en première position et tous les autres sont décalés d'une position vers la droite.

  \item La fonction ne renvoie rien (pas de \lstinline{return}) : elle renvoie implicitement \lstinline{None}. Le tableau est modifié \textbf{en place} (mutation) car les listes sont des objets mutables. C'est pourquoi \lstinline{scores} est modifié après l'appel, sans avoir besoin de récupérer une valeur de retour.
\end{enumerate}

Vérification : les résultats ont été confirmés par exécution.
\end{corrige}
